% !TeX program = pdflatex
% ------------------------------------------------------------------
% arXiv submission draft.  Suggested classes:
%   primary: math.AT ; secondary: cs.RO, cs.DC
% ------------------------------------------------------------------
\documentclass[12pt]{article}

\setlength{\oddsidemargin}{.5cm}
\setlength{\evensidemargin}{.5cm}
\setlength{\textwidth}{6in}
\setlength{\textheight}{8in}
\setlength{\emergencystretch}{1.5em}

\usepackage{amsmath,amssymb,amsthm}
\usepackage{array,longtable}
\usepackage{float}
\usepackage{algorithm}
\usepackage{algpseudocode}
\usepackage[font=small,labelfont=bf]{caption}
\usepackage{booktabs}
\usepackage{hyperref}
\hypersetup{colorlinks=true,linkcolor=blue,citecolor=blue,urlcolor=blue}

\newcommand{\doi}[1]{doi:\,\href{http://dx.doi.org/#1}{#1}}
\newcommand{\TC}{\mathrm{TC}}
\newcommand{\cat}{\mathrm{cat}}
\newcommand{\SC}{\mathrm{SC}}
\newcommand{\SCst}{\SC_{\mathrm{strict}}}
\newcommand{\SD}{\mathrm{SD}}

\newtheorem{theo}{Theorem}[section]
\newtheorem{lemma}[theo]{Lemma}
\newtheorem{corollary}[theo]{Corollary}
\newtheorem{proposition}[theo]{Proposition}
\newtheorem{definition}[theo]{Definition}
\newtheorem{example}[theo]{Example}
\newtheorem{remark}[theo]{Remark}
\newtheorem{conjecture}[theo]{Conjecture}

\title{Simplicial complexity of spheres:\\
a chain-first annealer for explicit covers of
boundary-of-simplex spheres}

\author{M. Perez\footnote{Escuela Superior de F\'isica
y Matem\'aticas del Instituto Polit\'ecnico Nacional, Edificio 9, Unidad
Profesional Adolfo L\'opez Mateos, 07300 Mexico City, Mexico.}, and
J. Flores\footnotemark[1]}

\date{}

\begin{document}
\maketitle

\begin{abstract}
Gonz\'alez's simplicial complexity $\SCst(K)$ discretizes Farber's
topological complexity: at subdivision level zero it is a finite
integer, upper-bounded by exhibiting an explicit cover of $K\times K$
by contiguity domains.  The randomized greedy pipeline of Ortiz, Lara,
Gonz\'alez and Borat constructs such covers, but its published reach
stops at $1$-dimensional complexes: the experiments end with the circle
and the wedge of two circles, the $2$-sphere was deferred to cluster
hardware, and in our faithful reimplementation the greedy search never
completed a single covering of $\partial\Delta^3\times\partial\Delta^3$
in $120$-hour jobs on $96$ cores.  We diagnose two mechanisms behind
this stall---an objective flat over the number of chain defects, the
very quantity the walk controls, and untracked chain entries drifting
to noise---and repair both with a \emph{chain-first annealer}: whole
systems of $k$ complete contiguity chains, evolved in parallel on a
single GPU under an energy that grades every uncovered facet by the
defect count of its best chain and penalises defects in all chains at
once.  The corrected search covers the spheres $S^1$ through $S^5$ by
explicit three-domain systems, in $22$ seconds ($S^2$), $153$ seconds
($S^3$), $13$ minutes ($S^4$) and $34$ minutes ($S^5$), certifying
$\SCst(\partial\Delta^{n+1})\le2$ for all $n\le5$, with the sharp
values $\SCst=\TC=2$ on $S^2$, on $S^4$ and on the cube boundary---a
triangulation in which product constructions require eight domains.  A
pattern in the search's own plateaus then becomes a theorem: every
rectangle $\operatorname{star}(v)\times\operatorname{star}(w)$ is a
contiguity domain, giving $\SCst(\partial\Delta^{n+1})\le2$ for
\emph{all} $n\ge1$ and $\SCst(\partial\Delta^{n+1})=\TC(S^n)=2$ at
level zero for every even $n$.  Every reported system is an audited
certificate, replayed link by link by an independent checker sharing no
code with the search; the full certificate set re-verifies in about two
minutes on CPU alone.
\end{abstract}

\medskip\noindent
\emph{Keywords and phrases:} topological complexity; simplicial
complexity; contiguity chains; piecewise linear motion planning;
simulated annealing; GPU computing.

\medskip\noindent
\emph{2020 Mathematics Subject Classification:} Primary: 55M30, 68W20.
Secondary: 68T40, 55U10, 90C27.

\section{Introduction}\label{sec:intro}

Farber's topological complexity $\TC(X)$ \cite{F03} measures the
smallest number of continuous rules needed to move a system between any
two configurations of $X$.  Gonz\'alez's simplicial complexity
\cite{G18} discretizes the question: replace continuous motion planners
by \emph{simplicial} ones, defined on pieces of the ordered product
$K\times K$ of a triangulation $K$ of $X$, where ``continuous'' becomes
``joined by a contiguity chain of simplicial maps''
(Section~\ref{sec:back}).  The value $\SCst(K)$ at subdivision level
zero is a finite, checkable integer---but it is only as real as the
best explicit cover of $K\times K$ by contiguity domains, and finding
such a cover is a search problem over chains of simplicial maps whose
cost had so far pushed every target past dimension one out of reach.

Ortiz, Lara, Gonz\'alez and Borat \cite{OLGB21} gave a pipeline of
randomized greedy algorithms for exactly this search: grow one domain
at a time from a seed facet by budgeted random walks, then merge and
prune the resulting cover (Section~\ref{sec:back}).  Their published
experiments draw its cost curve: the circle $\partial\Delta^2$ ($18$
maximal facets in the product) takes about a minute; the wedge of two
circles ($50$ facets) takes $25$ minutes at level~$0$; a single
barycentric subdivision of the circle takes about $10$ hours; and the
$2$-sphere $\partial\Delta^3\times\partial\Delta^3$, with its $96$
maximal facets, was explicitly deferred to cluster-scale hardware.

The cluster estimate proved optimistic.  A C++/MPI predecessor of the
present work ran the same algorithms on a $96$-core cluster in
$120$-hour jobs and never completed one covering of
$\partial\Delta^3\times\partial\Delta^3$: a single failed growth
attempt cost about $3$ seconds per candidate walk, domains froze after
a few dozen facets, and dozens of consecutive retry passes flipped
nothing at full computational cost.  A faithful single-GPU port of the
same pipeline---bit-for-bit against a CPU reference, on one Tesla
T4---did produce the first explicit cover of the $2$-sphere model, in
$4.4$ minutes, attaining $\SCst(\partial\Delta^3)=\TC(S^2)=2$; but
brute force alone had reached its ceiling, and on the $3$-sphere model
the best certified cover still had $20$ domains
($\SCst(\partial\Delta^4)\le19$).

This paper reports the method that broke that ceiling, and the theorem
the search had been trying to tell us.  There are four contributions,
in the order the reader meets them.

\begin{itemize}
\item \textbf{Diagnosis} (Section~\ref{sec:fail}).  The greedy walk
dies in two documented ways: its score is flat over the number of
chain defects---the quantity the walk actually controls---and chain
entries that are momentarily irrelevant receive no signal and drift to
noise.
\item \textbf{A chain-first annealer} (Section~\ref{sec:annealer}).
Instead of growing domains facet by facet, we anneal whole systems of
$k$ complete contiguity chains in parallel on one GPU, under an energy
that grades every uncovered facet by the defect count of its best
chain and penalises defects in all chains at once.  The search covers
$S^1$ through $S^5$ by explicit three-domain systems---$22$ seconds
for the $2$-sphere, $34$ minutes for the $5$-sphere---certifying
$\SCst(\partial\Delta^{n+1})\le2$ for all $n\le5$, sharply ($=\TC$) on
$S^2$, $S^4$ and the cube boundary (Table~\ref{tab:annealer}).
\item \textbf{The product-domains theorem}
(Section~\ref{sec:products}).  Every plateau of every two-chain search
sat at exactly the coverage attained by two product rectangles
$\operatorname{star}(v)\times\operatorname{star}(w)$.  We prove every
such rectangle is a contiguity domain (Theorem~\ref{thm:prod}), giving
$\SCst(K)\le\rho(K)-1$ for a rectangle-cover constant $\rho(K)$
computable by a small integer program (Corollary~\ref{cor:rho}); for
the simplex boundary $\rho(\partial\Delta^{n+1})=3$, hence
\[
\SCst(\partial\Delta^{n+1})\le 2\ \ \text{for all }n\ge1,
\qquad
\SCst(\partial\Delta^{n+1})=\TC(S^n)=2\ \ \text{for even }n\ge2,
\]
sharp at subdivision level $0$ (Corollary~\ref{cor:spheres}) and
including, to our knowledge, the first sharp level-$0$ value beyond
dimension $2$ in this framework ($S^4$).
\item \textbf{Certificates, not claims} (Section~\ref{sec:repro}).
Every reported system is replayed deterministically and re-audited by
an independent checker that rebuilds each triangulation from scratch
and shares no code with the search; the full certificate set verifies
in about two minutes on CPU alone, and source and certificates are
public.
\end{itemize}

The theorem and the annealer are complementary, and
Sections~\ref{sec:annealer} and \ref{sec:results} keep them in tension: on the cube boundary,
where $\rho=8$ caps every product cover at $\SCst\le7$, the annealer
still finds the optimum three domains; on the dodecahedron boundary a
$20$-hour ladder campaign finds no cover at all; and the exact
level-$0$ values on odd-dimensional spheres remain open
(Conjecture~\ref{conj}).  The results, old and new, are collected in
Table~\ref{tab:records}.

\section{Background and notation}\label{sec:back}

We keep the framework of \cite{OLGB21} and recall only what is used
later.  Throughout, $K$ is a finite connected simplicial complex on an
ordered vertex set $V$, and $K\times K$ is the ordered (``staircase'')
product triangulation of $\|K\|\times\|K\|$; its maximal facets lie
over pairs $(s,t)$ of facets of $K$.

Two simplicial maps $\varphi,\varphi':J\to K$ are \emph{$1$-contiguous}
if $\varphi(\sigma)\cup\varphi'(\sigma)$ is a simplex of $K$ for every
facet $\sigma$ of $J$; a \emph{contiguity chain} is a sequence
$\varphi=\varphi_0\sim\varphi_1\sim\cdots\sim\varphi_c=\varphi'$ of
consecutively $1$-contiguous maps.  A subcomplex $J\subseteq L$ is a
\emph{contiguity domain} for $\psi,\psi':L\to K$ if the restrictions
$\psi|_J,\psi'|_J$ are joined by a contiguity chain, and the strict
contiguity distance $\SD(\psi,\psi')$ is one less than the minimum
number of contiguity domains covering $L$.  For the two projections
$\pi_1,\pi_2:K\times K\to K$,
\[
\SCst(K)=\SD(\pi_1,\pi_2)\ge \TC(\|K\|),
\]
with equality in the limit of barycentric subdivisions
\cite{OLGB21,G18}.  A cover of $K\times K$ by $d$ contiguity domains is
a \emph{system of $d$ piecewise-linear motion planners} and certifies
$\SCst(K)\le d-1$.  All statements in this paper are at subdivision
level $0$, and we read the facets of $K\times K$ as a covering problem
over the \emph{facet pairs} $(s,t)$.  When $K$ is pure of dimension
$d$, every pair carries the same number $\binom{2d}{d}$ of staircase
facets, and a facet count is a pair count up to that factor; when $K$
is not pure, the number of staircase facets over $(s,t)$ depends on
$\dim s$ and $\dim t$.  In either case a coverage figure counted in
facets becomes a count of covered \emph{pairs} only after checking
that each pair is covered in its entirety, since a domain may contain
some of the staircase facets over a pair without containing them all.

The pipeline of \cite{OLGB21} grows one domain at a time: starting from
a seed facet it repeatedly calls \textsc{LocalSearch}---a random walk
with budget $M$ that tries to extend a contiguity chain over one more
facet---then merges and prunes the resulting cover.  Its cost profile,
and the way it fails, is the subject of the next section.

\section{How the greedy dies}\label{sec:fail}

Three autopsy findings, each verified experimentally.

\textbf{Flat energy.}  A \textsc{LocalSearch} walk is scored by whether
the \emph{whole} chain works, so a candidate whose chain has one bad
link scores essentially like a candidate with forty bad links: there is
no gradient over the number of defects, the quantity the walk actually
controls.  We observed this regime directly in an early version of the
chain-first annealer: with a defect penalty of $\beta=0.05$ it never
solves even the circle (the smallest nontrivial target, $18$ facets),
stalling run after run at $13$--$15$ covered facets, indefinitely.
That is a documented failure mode of our own scoring function; that the
same flatness is what stalls the pipeline of \cite{OLGB21} is a reading
of the two behaviours, not a causal link we have isolated.

\textbf{Chain drift.}  When several chains compete to cover the same
facets and the score takes the minimum over chains, entries of a chain
that are momentarily irrelevant receive no signal at all and drift to
noise; when a facet should migrate from one chain to another, the
receiving chain has been randomized exactly where it matters.  We read
this as the likely microscopic mechanism of the macroscopic freeze
observed from the beginning---on the $96$-core cluster, growth froze
after the first few dozen facets of each domain, and retry passes
flipped nothing---though we have not established that link
experimentally.  The
wedge-of-two-circles reproduction of the first phase of this project
exhibited the same
signature at small scale: exhaustive pairwise-merge probing showed the
greedy's recorded partition to be a genuine local optimum of its own
dynamics---stable, and one domain above the true optimum.

\textbf{Quadratic accounting.}  Each step of the original walk rescans
the whole domain and re-runs a shortest-path computation from scratch:
about $70\,\mu$s per step in the C++/MPI implementation, so one failed
growth attempt at budget $M=40000$ costs about $3$ seconds per
candidate, sequentially per core.  $96$ cores times $120$ hours of this
is a large number of stalls, not a large covering.

The repair of the accounting is engineering (batched walks, incremental
step conditions, precomputed extension sets; on one Tesla T4 the
batched kernel issues about $10^{10}$ walk steps in $4.4$ minutes).
The repair of the two diseases is the energy of the next section, and
it is what made the corrected search worth running.

\section{The chain-first annealer}\label{sec:annealer}

Instead of growing domains facet by facet, we anneal \emph{whole
candidate systems}.  A state is a tuple of $k$ chains.  A chain is a
sequence $C=(\varphi_0,\varphi_1,\dots,\varphi_{c_{\max}})$ of maps of
vertex sets with fixed endpoints $\varphi_0=\pi_1$ and
$\varphi_{c_{\max}}=\pi_2$; the interior entries are arbitrary vertex
assignments and are \emph{not} required to be simplicial maps
$K\times K\to K$.  In the implementation $c_{\max}$ counts links, so a
stored chain holds $c_{\max}+1$ maps.

A facet $\sigma$ of $K\times K$ is \emph{covered} by $C$ when every
link of $C$ satisfies the $1$-contiguity condition on $\sigma$, and the
\emph{domain} $\mathrm{Dom}(C)$ is the subcomplex generated by the
facets that $C$ covers.  Restricted to $\mathrm{Dom}(C)$ the entries
are simplicial and consecutive entries are $1$-contiguous, so
$C|_{\mathrm{Dom}(C)}$ is a contiguity chain from
$\pi_1|_{\mathrm{Dom}(C)}$ to $\pi_2|_{\mathrm{Dom}(C)}$ in the sense
of Section~\ref{sec:back}.  It is this restriction, not the ambient
vertex assignment, that a certificate reports.

For a facet $\sigma$ and the $j$-th chain of a state, let $b_j(\sigma)$
be the number of links of that chain which violate the contiguity
condition on $\sigma$, and let $m(\sigma)=\min_j b_j(\sigma)$ be the
defect count of the best chain for that facet.  With $\mathrm{pen}(0)=0$
and $\mathrm{pen}(b)=1+\beta b$ for $b>0$, the energy of a state is
\[
E \;=\; \sum_{\sigma} w_\sigma\,\mathrm{pen}\bigl(m(\sigma)\bigr)
\;+\; \gamma\sum_{j}\sum_{\sigma} b_j(\sigma).
\]
The facet weights are $w_\sigma=1$ in the unguided runs; a guided
restart may raise the weight of facets that an earlier stage left
uncovered.  The first sum is the objective: it vanishes exactly on the
facets some chain already covers, and $\beta$ grades the rest by how
far the \emph{best} chain is from covering them, which is the gradient
a flat score lacks.  The second sum penalises defects in \emph{all}
chains, not only the best one; it is the cleanliness term that
suppresses drift---without it ($\gamma=0$) long chains fail again, and
$\gamma=0.1$ is measurably too weak.  \emph{Coverage}, the quantity
reported in the tables, is the number of facets with $m(\sigma)=0$; it
is not itself a term of $E$.

The word \emph{energy} is the standard simulated-annealing term
\cite{KGV83,Cerny85}, not a physical quantity: in the Metropolis
dynamics \cite{Metropolis53}, the cost function of a combinatorial
search plays the role that energy plays in statistical mechanics, and
the temperature $T$ entering the acceptance rule $e^{-\Delta E/T}$
below is a control parameter measured in the same (dimensionless)
units as $E$.

A population of $S$ states anneals in parallel on the GPU under a
Metropolis kernel of its own, distinct from the \textsc{LocalSearch}
kernel of the first phase (Section~\ref{sec:intro}).  The population is
not drawn with independent uniform interior maps: each state is built
by sweeps that flip entries of $\pi_1$ towards $\pi_2$ in random order,
so that every initial chain already interpolates between the two
projections.  A proposal chooses one chain and one vertex uniformly at random and
changes the value assigned to that vertex, either at a single interior
position or over a block of two to eight consecutive positions; the
block proposals are what we call \emph{track moves}, and they are drawn
with probability $p_{\mathrm{track}}$.  The new value is copied from a
map adjacent to the touched block with probability $1/2$, and otherwise
drawn uniformly from the remaining vertices; copy moves let a locally
good assignment propagate along the chain.  Track moves are a move
shape, not an explicit bookkeeping of the facets a chain has been
assigned.  The temperature follows a geometric schedule \emph{within}
each epoch, and the incumbent best state is kept by highest unweighted
coverage, which need not be the state of lowest energy.  Every second
epoch the population undergoes selection and replacement---the weakest
quarter of the states is overwritten by copies of the elite and a few
slots are re-initialised at random---so the $S$ trajectories are not
independent for the whole run.  Once the incumbent covers every facet
the run stops, and the incumbent's chains are used to \emph{partition}
the facets (each facet goes to the first chain covering it), each
restricted chain is shortened by the greedy farthest-jump reduction of
\cite[Algorithm 2]{OLGB21}, and the resulting system is verified map by
map, link by link, facet by facet.  The output is a certificate, not a
claim (Section~\ref{sec:repro}).
Algorithm~\ref{alg:anneal} gives the full loop in the style of
\cite{OLGB21}.

\begin{algorithm}[H]
\caption{\textsc{ChainAnneal}($K$, $k$, $c_{\max}$; $S$, $\beta$,
$\gamma$, $p_{\mathrm{track}}$): anneal whole systems of $k$ chains over $K\times K$.}
\label{alg:anneal}
\begin{algorithmic}[1]
\Require a \emph{state} is $(C_1,\dots,C_k)$ with
$C_i=(\varphi_{i,0},\dots,\varphi_{i,c_{\max}})$,
$\varphi_{i,0}=\pi_1$, $\varphi_{i,c_{\max}}=\pi_2$, interior entries
arbitrary vertex assignments; population size $S$
\State initialise $S$ states by random-order sweeps that flip entries
of $\pi_1$ towards $\pi_2$
\State $\mathcal{C}^*\gets$ incumbent (highest coverage seen so far)
\For{epoch $=1,2,\dots$}
\ForAll{states $\mathcal{C}$ of the population, in parallel on the GPU}
\State perform $M$ Metropolis steps, the temperature $T$ following a
geometric schedule $T_0\to T_1$ \emph{inside} the epoch: pick one chain
$j$ and one vertex $w$ uniformly and propose a new value $u$ for
$\varphi_{j,i}(w)$, at a single interior position $i$ or (with
probability $p_{\mathrm{track}}$) over a block of two to eight
consecutive positions $i_1,\dots,i_2$; with probability $1/2$ the value
$u$ is copied from a neighbouring map, $\varphi_{j,i_1-1}(w)$ or
$\varphi_{j,i_2+1}(w)$, otherwise drawn uniformly from the remaining
vertices; accept by $e^{-\Delta E/T}$
\Statex \hspace{2em} with
$E=\sum_\sigma w_\sigma\,\mathrm{pen}(m(\sigma))
+\gamma\sum_j\sum_\sigma b_j(\sigma)$
\EndFor
\State update $\mathcal{C}^*$ by highest unweighted coverage
\Comment{not necessarily lowest $E$}
\If{$\mathcal{C}^*$ covers every facet of $K\times K$}
\State partition: each facet goes to the first chain of $\mathcal{C}^*$
covering it \Comment{facet $\mapsto$ domain}
\State shorten each restricted chain by the greedy farthest-jump
reduction \cite[Algorithm~2]{OLGB21}
\State verify: every map simplicial and every link $1$-contiguous on
every facet of its domain, chain from $\pi_1|_J$ to $\pi_2|_J$
\State \Return $(\text{partition},\text{chains})$
\Comment{a certificate for $\SCst(K)\le k-1$}
\EndIf
\State every second epoch, overwrite the weakest quarter of the
population with copies of the elite and re-initialise a few slots at
random \Comment{selection and replacement}
\EndFor
\end{algorithmic}
\end{algorithm}

Table~\ref{tab:annealer} summarizes the runs and Table~\ref{tab:params}
lists the search parameters of each one.  The plateau rows are as
instructive as the successes, and we return to them in
Section~\ref{sec:products}: the two-chain runs on $\partial\Delta^3$,
$\partial\Delta^4$ and $\partial\Delta^5$ stop at coverage values that
are exactly those attained by suitable two-rectangle constructions.

\begin{table}[H]
\centering\footnotesize
\setlength{\tabcolsep}{3.5pt}
\begin{tabular}{@{}llr@{\hspace{5pt}}lrlp{3.0cm}@{}}
\toprule
target & triangulation $K$ & facets & predicted & best coverage & wall time & outcome (domains found) \\
\midrule
$S^1$ & $\partial\Delta^2$ & $18$ & $2$ & $18/18$ & $<15$\,s
  & $2$ domains, verified; $\SCst{\le}2$ ($=1$ known \cite{OLGB21})\\
$S^2$ & $\partial\Delta^3$ & $96$ & $3$ & $96/96$ (epoch $1$) & $22$\,s
  & $54+29+13$; $\SCst{=}2{=}\TC$\\
$S^3$ & $\partial\Delta^4$ & $500$ & $3$ & $500/500$ (epoch $1$) & $153$\,s
  & $309+130+61$; $\SCst{\le}2$\\
$S^4$ & $\partial\Delta^5$ & $2520$ & $3$ & $2520/2520$ (epoch $1$) & $13$\,min
  & $1634+653+233$; $\SCst{=}2{=}\TC$\\
$S^5$ & $\partial\Delta^6$ & $12348$ & $3$ & $12348/12348$ (epoch $2$) & $34$\,min
  & $7061+4239+1048$; $\SCst{\le}2$\\
\midrule
$S^1\vee S^1$ & $S^1\vee S^1$ & $50$ & $3$ & $50/50$ (epoch $25$) & $41$\,s
  & $18+16+16$; $\SCst{=}2$\\
$S^2$ & cube boundary & $864$ & $8$ ($\rho$) & $864/864$ (epoch $22$) & ${\approx}24$\,min$^\dagger$
  & $482+236+146$; $\SCst{=}2{=}\TC$\\
\midrule
\multicolumn{7}{@{}l}{\emph{Two-chain plateau runs ($k=2$): the search seeks a cover by $2$ domains and stalls.}}\\
$S^2$ & $\partial\Delta^3$ & $96$ & $2$ & $84/96$ & (plateau)
  & $84/96=14/16$ of facet pairs\\
$S^3$ & $\partial\Delta^4$ & $500$ & $2$ & $460/500$ & (plateau)
  & $460/500=23/25$ of facet pairs, all seeds\\
$S^4$ & $\partial\Delta^5$ & $2520$ & $2$ & $2380/2520$ (epoch $1$) & $12$\,min
  & $2380/2520=34/36$ of facet pairs, as predicted\\
$S^2$ & cube boundary & $864$ & $2$ & --- & ---
  & $\TC(S^2)=2$: no $2$-domain cover exists\\
\bottomrule
\end{tabular}
\caption{Chain-first annealer, one Tesla T4.  The \emph{predicted}
column is the number of domains the run seeks: $3$ for the $k=3$ runs
(which the product-domains theorem, Section~\ref{sec:products},
guarantees suffices on the simplex boundaries) and $2$ for the plateau
runs; on the cube boundary the product bound is $\rho=8$, so the three
domains found beat every product cover.  Domain sizes are those of the
certified cover.  The per-run values of $c_{\max}$, the population size
and the move budget are \emph{not} uniform across rows; they are listed
in Table~\ref{tab:params}.  The plateaus reported here were stable
across the seeds and the continuations we ran.  $^\dagger$The cube
entry is a three-rung ladder totalling about $24$ minutes---cold
rung ($303$\,s, $662/864$), warm rung ($814$\,s, $820/864$), guided
warm rung ($304$\,s, complete); rungs, warm starts and guided weights
are defined below.  The $304$\,s quoted above is the final guided rung
alone, not the whole sequence.}
\label{tab:annealer}
\end{table}

\begin{table}[H]
\centering\footnotesize
\begin{tabular}{@{}lrrrr@{}}
\toprule
run & $c_{\max}$ & $S$ & moves/epoch & seed \\
\midrule
$\partial\Delta^3$ (S$^2$) & $20$ & $4096$ & $20000$ & $51$\\
$\partial\Delta^4$ (S$^3$) & $40$ & $4096$ & $20000$ & $52$\\
$\partial\Delta^5$ (S$^4$) & $40$ & $4096$ & $20000$ & $61$\\
$\partial\Delta^6$ (S$^5$) & $16$ & $4096$ & $5000$ & $71$\\
wedge $S^1\vee S^1$ & $16$ & $4096$ & $5000$ & $71$\\
cube boundary (final guided stage) & $16$ & $4096$ & $5000$ & $84$\\
\bottomrule
\end{tabular}
\caption{Search parameters of the runs reported in
Table~\ref{tab:annealer}; all six used $\beta=5$, $\gamma=0.5$ and
$p_{\mathrm{track}}=0.3$.  These settings should not be read back onto
the circle runs, which used different values.}
\label{tab:params}
\end{table}

\emph{Ladders of restarts: rungs, warm starts, guided weights.}  The
runs of Table~\ref{tab:annealer} were single calls to
Algorithm~\ref{alg:anneal}, with one exception.  The cube boundary
certificate was reached by a short \emph{ladder}: a sequence of bounded
calls, each of which we call a \emph{rung}, chained so that each rung
starts where the previous one stopped.  A \emph{cold} rung is an
ordinary call to Algorithm~\ref{alg:anneal}; on the cube it covered
$662$ of $864$ facets in $303$\,s.  A \emph{warm} rung initialises its
population from the previous rung's incumbent: half of the population
receives the $k$ incumbent chains, each shortened on its domain by the
greedy reduction and padded with copies of $\pi_2$ up to $c_{\max}$
links, and the other half is built from fresh sweeps as before; on the
cube it continued to $820/864$ in $814$\,s.  A \emph{guided} rung
additionally rebalances the energy during the search: after each epoch
every weight decays back towards $1$ by $w_\sigma\gets
1+0.7\,(w_\sigma-1)$, and every facet the incumbent still misses
receives an additive increment $g$, capped at $g_{\max}$; reported
coverage stays unweighted.  A guided warm rung closed the remaining
$44$ facets in $304$\,s.  Guidance without the seed is not enough: a
guided \emph{cold} rung stalls at $858/864$ with the whole population
pinned at $856$--$858$---the guide needs a good incumbent, not only a
budget.

Why short rungs at all?  The coverage gained within a single warm rung
is \emph{front-loaded}: measured on the dodecahedron campaign below, it
fits a sublinear law $\Delta\mathrm{cov}(t)\approx c\,t^{0.8}$ (one
rung gained $+33$ facets in its first $7$ minutes, $+58$ in the first
$15$, and $+170$ in the full hour).  Early in a rung the population
claims the facets accessible from the incumbent's chain topology; once
those are exhausted, further time buys rapidly diminishing returns.
Every rung also pays a fixed overhead $\tau$ (kernel compilation and
host--device transfer, $\tau\approx40$\,s on the machine used here), so
the coverage harvested per GPU-hour, $c\,t^{0.8}/(t+\tau)$, peaks at
$T^{*}_{\mathrm{rung}}=\tfrac{0.8}{1-0.8}\,\tau=4\tau\approx2.7$\,min.
Algorithm~\ref{alg:ladder} abstracts the pattern.

\begin{algorithm}[H]
\caption{\textsc{WarmLadder}($K$, $k$, $c_{\max}$; $S$, $\beta$,
$\gamma$, $p_{\mathrm{track}}$; per-rung budget $T_{\mathrm{rung}}$,
guided increment $g$, cap $g_{\max}$, max rungs $R_{\max}$).
Algorithm~\ref{alg:anneal} run as a ladder: warm seeding between rungs,
guided reweighting within each rung.}
\label{alg:ladder}
\begin{algorithmic}[1]
\State run \textsc{ChainAnneal}$(K,k,c_{\max};S,\beta,\gamma,
p_{\mathrm{track}})$ for at most $T_{\mathrm{rung}}$ (\emph{cold} rung,
$w_\sigma\equiv1$); set $\mathcal{C}^{*}\gets$ best state found
\For{$r=1,2,\ldots,R_{\max}$}
  \If{$\mathcal{C}^{*}$ covers all facets of $K\times K$}
    \State\Return $\mathcal{C}^{*}$
    \Comment{certificate: $\SCst(K)\le k-1$}
  \EndIf
  \State \emph{warm seeding}: half the population $\gets$ the chains of
  $\mathcal{C}^{*}$, each reduced on its domain and padded with $\pi_2$
  up to $c_{\max}$ links; the other half is freshly swept
  \State \emph{guided rung}: run \textsc{ChainAnneal} for at most
  $T_{\mathrm{rung}}$ with adaptive weights---after each epoch,
  $w_\sigma\gets1+0.7\,(w_\sigma-1)$ for every facet $\sigma$ and
  $w_\sigma\gets\min(w_\sigma+g,\,g_{\max})$ for every facet the
  incumbent does not cover; update $\mathcal{C}^{*}\gets$ best state
  found \Comment{coverage reporting stays unweighted}
\EndFor
\State\Return $\mathcal{C}^{*}$
\Comment{partial cover; no certificate}
\end{algorithmic}
\end{algorithm}

To probe how far ladders extend the reach of the search, we applied
Algorithm~\ref{alg:ladder} to the dodecahedron boundary: a fan
triangulation of $S^2$ with $20$ vertices and $36$ triangles, yielding
$K\times K$ with $7776$ product facets.  Its simplex density is $3.2\%$
(compared to $21.4\%$ for the cube boundary), and the dual graph of $K$
has diameter $7$ against $4$ for the cube---both figures predict
substantially harder search.  We ran five long warm escalations of up
to one hour each, then a systematic $2.5$-minute ladder of $172$ short
rungs executed unattended over $12$ GPU-hours.
Table~\ref{tab:dodeca} records the complete campaign.

\begin{table}[H]
\centering\footnotesize
\begin{tabular}{@{}lrrrr@{}}
\toprule
stage & time & coverage & gain & \% \\
\midrule
cold start ($c_{\max}=40$, $S=4096$) &
  $60$\,min & $0\to 644$ & $+644$ & $8.3$ \\
warm escalation 1 &
  $60$\,min & $644\to 1208$ & $+564$ & $15.5$ \\
warm escalation 2 &
  $60$\,min & $1208\to 1579$ & $+371$ & $20.3$ \\
warm escalation 3 &
  $60$\,min & $1579\to 1749$ & $+170$ & $22.5$ \\
warm escalation 4 &
  $15$\,min & $1749\to 1802$ & $+53$ & $23.2$ \\
warm escalation 5 &
  $7$\,min & $1802\to 1831$ & $+29$ & $23.5$ \\
\midrule
2.5-min ladder, R.\,1--20 &
  $55$\,min & $1831\to 2064$ & $+233$ & $26.5$ \\
2.5-min ladder, R.\,22--41 &
  $58$\,min & $2075\to 2218$ & $+143$ & $28.5$ \\
2.5-min ladder, R.\,42--61 &
  $58$\,min & $2218\to 2310$ & $+92$ & $29.7$ \\
2.5-min ladder, R.\,62--81 &
  $58$\,min & $2310\to 2400$ & $+90$ & $30.9$ \\
2.5-min ladder, R.\,82--101 &
  $58$\,min & $2400\to 2441$ & $+41$ & $31.4$ \\
2.5-min ladder, R.\,102--121 &
  $58$\,min & $2441\to 2487$ & $+46$ & $32.0$ \\
2.5-min ladder, R.\,122--141 &
  $58$\,min & $2487\to 2523$ & $+36$ & $32.4$ \\
2.5-min ladder, R.\,142--172 + probes &
  ${\approx}3.5$\,h & $2523\to 2602$ & $+79$ & $33.5$ \\
\midrule
\textbf{total} &
  ${\approx}20$\,h & $0\to 2602$ & $+2602$ & $33.5$ \\
\bottomrule
\end{tabular}
\caption{Complete warm-start rung-ladder campaign on the dodecahedron boundary
($7776$ product facets of $K\times K$, $k=3$).
Upper block: five long warm escalations ($c_{\max}=40$, $g=3$, $g_{\max}=11$),
each reseeded from the previous incumbent.
Lower block: $2.5$-minute ladder run unattended over $12$\,GPU-hours,
summarised in batches of $20$ rungs; the final row combines the last partial
batch with extended window probes and a $70$-trial basin survey.
Mean gain per rung falls from $11.7$ in the first batch to below $1$ in
the last; no complete cover was found.
The obstacle is the sparsity of $K$: a dual-graph diameter of $7$ requires
coordinated walks that grow geometrically harder to satisfy simultaneously
as coverage rises.}
\label{tab:dodeca}
\end{table}

We close the section with four observations on
Table~\ref{tab:annealer}.

\emph{Certified values, not just bounds.}  The successes pin the
invariant itself: for $S^2$ and $S^4$ the covers found, together with
the lower bound $\TC=2$ of Section~\ref{sec:products}, force
$\SCst=\TC=2$, so the search independently reproduces at level $0$ the
optimum that theorem proves in general.

\emph{Near-product shape.}  The winning domains on the simplex
boundaries are close to star rectangles, but measurably not equal to
them: on $\partial\Delta^5$ the chains span $30$, $28$ and more facet
pairs where a rectangle spans $25$, with $62$, $3$ and $67$ facets
lying outside the closest rectangle, and on $\partial\Delta^4$ the
largest single domain holds $329$ facets where a rectangle holds $320$.
The annealer spends a small, quantifiable amount of non-product
freedom.

\emph{The plateau value is structural; the plateau shape is not.}
The stalled two-chain runs stop exactly at the coverage of the best
two-rectangle constructions.  This does not mean the search found
rectangles: among the saved states, one $\partial\Delta^4$ run ends in
two genuine rectangles of $320$ facets each, union $460$, while two
other runs reach the same $460$ with domains contained in no rectangle
$\operatorname{star}(v)\times\operatorname{star}(w)$ at all, and the
$\partial\Delta^5$ run at $2380/2520$ is likewise not a union of two
rectangles.\footnote{The runs cited are \texttt{s3\_k2\_v2\_track}
(two rectangles of $320$), \texttt{s3\_k2\_A} (domains of $284$ and
$265$, the second in no rectangle) and \texttt{s3\_k2\_v2\_point}
(domains of $297$ and $317$, neither in a rectangle; union still
$460$) on $\partial\Delta^4$, and \texttt{s4\_k2\_v2} ($2380/2520$)
on $\partial\Delta^5$.  The saved states accompany the certificates of
Section~\ref{sec:repro}.}  What is stable across seeds is the coverage
value, not the domains that realise it---and explaining \emph{that
number} is what makes the next section possible.

\emph{Non-product freedom pays.}  The cube boundary, an $8$-vertex,
$12$-face triangulation of $S^2$ distinct from $\partial\Delta^3$,
shows the extreme case: the star-rectangle bound $\rho=8$ certifies
only $\SCst\le7$, while the free-form three-chain search reaches
$\SCst=2=\TC(S^2)$---five domains below the best certificate any
product construction could give.

\section{The product domains theorem}\label{sec:products}

Write $\operatorname{star}(v)$ for the set of facets of $K$ containing
the vertex $v$, and $c_x$ for the constant map to the vertex $x$.

\begin{theo}[product domains]\label{thm:prod}
For vertices $v,w$ of $K$, let $D(v,w)\subseteq K\times K$ be the
subcomplex generated by the maximal facets lying over pairs $(s,t)$ of
facets of $K$ with $v\in s$ and $w\in t$.  Then $D(v,w)$ is a
contiguity domain for
$(\pi_1,\pi_2)$: if $v=v_0,v_1,\dots,v_m=w$ is a path in the
$1$-skeleton of $K$, then on $D(v,w)$
\[
\pi_1 \;\sim\; c_{v_0}\;\sim\; c_{v_1}\;\sim\;\cdots\;\sim\;
c_{v_m}\;\sim\; \pi_2,
\]
a chain of $m+2$ links ($2$ links when $v=w$).
\end{theo}

\begin{proof}
The generating maximal facets are maximal in $K\times K$, so none is a
face of another; the facets of the subcomplex they generate are exactly
the generators, and since $1$-contiguity is a condition on facets
(Section~\ref{sec:back}), it suffices to check the chain on those.  A
facet $\sigma\in D(v,w)$ lying over $(s,t)$ has $\pi_1(\sigma)=s$,
$\pi_2(\sigma)=t$ (a staircase path visits every row and every column,
so all vertices of $s$ and of $t$ occur), and $c_x(\sigma)=\{x\}$, so
every map in the chain is simplicial on $\sigma$.  Consecutive maps are
$1$-contiguous on $\sigma$: $\pi_1(\sigma)\cup c_v(\sigma)=s\cup\{v\}=s$
since $v\in s$; $c_{v_i}(\sigma)\cup c_{v_{i+1}}(\sigma)=\{v_i,v_{i+1}\}$
is an edge of $K$; and $c_w(\sigma)\cup\pi_2(\sigma)=t$.
\end{proof}

We emphasize the provenance: this statement was not guessed, it was
\emph{found}.  The plateaus of Table~\ref{tab:annealer} sit at exactly
the coverage two such rectangles achieve---though, as recorded in
Section~\ref{sec:annealer}, the domains realising that coverage need not
themselves be rectangles---and in naming the value the theorem converts
a search problem into a covering problem.

\begin{corollary}[a search-free upper bound]\label{cor:rho}
Let $\rho(K)$ be the minimum number of rectangles $D(v,w)$ covering all
facet pairs of $K\times K$, i.e.\ the minimum size of a set
$R\subseteq V\times V$ such that for every pair $(s,t)$ of facets of
$K$ there is $(v,w)\in R$ with $v\in s$, $w\in t$.  Then
\[
\SCst(K)\;\le\;\rho(K)-1.
\]
\end{corollary}

\begin{proof}
Every facet pair $(s,t)$ has some $(v,w)\in R$ with $v\in s$, $w\in
t$, so the domains $D(v,w)$, $(v,w)\in R$, cover $K\times K$ (covers
need not be disjoint); Theorem~\ref{thm:prod} supplies the chains, and
$\rho(K)$ domains certify $\SCst(K)\le\rho(K)-1$.
\end{proof}

The number $\rho(K)$ is a set-cover optimum with $|V|^2$ candidate sets
over $|F|^2$ elements; a small integer program (HiGHS, a few hundred
constraints) solves it in seconds for every complex we consider.

\begin{corollary}[all boundary-of-simplex spheres]\label{cor:spheres}
For every $n\ge1$, $\rho(\partial\Delta^{n+1})=3$; hence
\[
\SCst(\partial\Delta^{n+1})\le 2\ \ \text{for all }n\ge1,
\quad
\SCst(\partial\Delta^{n+1})=\TC(S^n)=2\ \ \text{for even }n\ge2,
\]
sharp at subdivision level $0$.
\end{corollary}

\begin{proof}
The facets of $\partial\Delta^{n+1}$ are $s_a=V\setminus\{a\}$,
$a\in\{0,\dots,n+1\}$.  Take $R=\{(0,0),(1,1),(2,2)\}$: the pair
$(s_a,s_b)$ is missed by $(v,v)$ iff $v\in\{a,b\}$, so it is missed by
all three rectangles iff $\{0,1,2\}\subseteq\{a,b\}$, impossible.  Two
rectangles $(v,w),(v',w')$ never suffice: $(s_v,s_{w'})$ is missed by
both.  Hence $\rho=3$, and for even $n$ Farber's $\TC(S^n)=2$
\cite{F03} with $\TC(\|K\|)\le\SCst(K)$ \cite{G18} gives equality.
\end{proof}

\begin{remark}
The rectangle bound is not always sharp.  $\SCst(\partial\Delta^2)=1$
by the explicit two-domain cover of \cite{OLGB21}, while
Corollary~\ref{cor:rho} gives $2$; odd-dimensional spheres lie in
$\{1,2\}$ and the gap must come from non-product domains, as it does
for the circle.  The annealed two-chain searches on
$\partial\Delta^4$, however, never exceeded the coverage $460/500$ that
two rectangles already achieve, whatever the shape of the domains they
returned (Conjecture~\ref{conj}).  The cube boundary provides the starkest
known gap: $\rho(\text{cube})=8$ (verified by integer program,
Section~\ref{sec:results}), certifying only $\SCst\le 7$, while the
annealer finds $\SCst=2=\TC(S^2)$; product domains overshoot by a
factor of nearly four in the domain count.
\end{remark}

\section{Results}\label{sec:results}

Table~\ref{tab:records} places the new results alongside the best
previously known bounds.  The ``before'' column accumulates the
published pipeline \cite{OLGB21} and its reimplementations in the first
phase of this project (Section~\ref{sec:intro}): the $96$-core C++/MPI
version and the faithful single-GPU port.  The ``after'' column is
the chain-first annealer of Section~\ref{sec:annealer}.

\begin{remark}
The annealer certifies upper bounds, not optima.  A successful $k$-chain
run proves $\SCst(K)\le k-1$ but does not prove that $k-1$ is the
minimum; we are improving on the bounds achievable by the algorithms of
\cite{OLGB21}, not claiming that every reported value is a global
minimum.  The sharp equalities in the tables arise from independent
lower bounds---Farber's topological complexity for even-dimensional
spheres \cite{F03}, and prior constructions for the circle and the
wedge \cite{OLGB21}---matched by our certificates.  For spaces not yet
pinned by a matching lower bound the reported values are the best known
upper bounds.
\end{remark}

\begin{table}[H]
\centering\footnotesize
\begin{tabular}{@{}p{2.7cm}p{3.5cm}p{2.8cm}p{4.2cm}@{}}
\toprule
target & before & after & certificate \\
\midrule
circle $\partial\Delta^2$ ($18$)
  & $57$\,s \cite{OLGB21}; level $1\approx10$\,h
  & $<15$\,s
  & $\SCst\le 2$ ($=1$ known), audited\\
wedge $K_4\setminus e$ ($50$)
  & $25$\,min \cite{OLGB21}, $3$ domains
  & $41$\,s
  & $18+16+16$, $\SCst=2$, audited\\
cube boundary ($864$)
  & $\SCst\le 7$ (product bound, $\rho=8$)
  & $\approx 24$\,min total; $304$\,s final guided stage
  & $482+236+146$, $\SCst=2=\TC$, audited\\
$\partial\Delta^3$ (S$^2$, $96$)
  & never on $96$ cores, $120$\,h jobs;
    $4.4$\,min on T4 (first phase)
  & $22$\,s
  & $54+29+13$, $\SCst=2=\TC$, audited\\
$\partial\Delta^4$ (S$^3$, $500$)
  & $\SCst\le 19$ ($20$ domains, $\approx3$\,h GPU, first phase)
  & $153$\,s, one epoch
  & $309+130+61$, $\SCst\le 2$, audited\\
$\partial\Delta^5$ (S$^4$, $2520$)
  & not attempted
  & $13$\,min, one epoch
  & $1634+653+233$, $\SCst=2=\TC$, audited\\
$\partial\Delta^6$ (S$^5$, $12348$)
  & not attempted
  & $34$\,min, epoch $2$
  & $7061+4239+1048$, $\SCst\le 2$, audited\\
all $\partial\Delta^{n+1}$
  & S$^1$, S$^2$, S$^3$ case by case
  & $\SCst\le2$ for all $n$; $=\TC$ for even $n$
  & Corollary~\ref{cor:spheres}\\
\bottomrule
\end{tabular}
\caption{Then and now.  Facet counts of the verified partition are listed
in the certificate column.  Every entry was replayed and re-audited
independently of the search code.}
\label{tab:records}
\end{table}

Several entries deserve emphasis.

\emph{The wedge and the cube separate free-form search from product
covers.}  On the wedge of two circles the product bound gives only
$\SCst\le 3$ ($\rho=4$), while the annealer reproduces the sharp
certificate of \cite{OLGB21} in $41$\,seconds---finding a different
partition ($18+16+16$ vs.\ $20+18+12$) with shorter chains ($2$, $5$,
$5$ links vs.\ $8$, $11$, $10$).  On the cube boundary the gap is
more dramatic: $\rho=8$ forces any product cover to use at least $8$
domains, certifying only $\SCst\le7$, yet the chain-first annealer
covers all $864$ facets with just $3$ domains and confirms
$\SCst=2=\TC(S^2)$.  Neither result is reachable by Theorem~\ref{thm:prod}.

\emph{The $4$-sphere and $5$-sphere were priced out of the old world.}
$\partial\Delta^5\times\partial\Delta^5$ has $2520$ maximal facets and
$\partial\Delta^6\times\partial\Delta^6$ has $12348$; neither was
attempted with the greedy pipeline.  The chain-first annealer solves
them in $13$ and $34$~minutes, respectively, and delivers independently
audited certificates.  Level-$0$ sharpness
$\SCst(\partial\Delta^5)=\TC(S^4)=2$ is, to our knowledge, the first
sharp result beyond dimension $2$ in this framework.

\emph{Prediction became quantitative.}  Before the corrected runs,
rectangle arithmetic predicted two-chain plateaus at $460=23\cdot 20$
facets out of $500$ on $\partial\Delta^4$ and at $2380=34\cdot 70$ out
of $2520$ on $\partial\Delta^5$.  Both coverage values were observed
exactly, the latter in a single epoch.  The agreement is at the level of
the numbers: the domains returned at those values are often not
rectangles (Section~\ref{sec:annealer}), so what the plateau measures is
how far local dynamics reach, not that a two-chain optimum has product
structure.

\emph{The search and the theorem certify the same optimum.}  The
$2$-sphere and $4$-sphere values match $\TC(S^2)=\TC(S^4)=2$ given by
Farber's theorem \cite{F03}, independently certifying sharpness at
subdivision level~$0$.


\section{Certificates, auditing, reproduction}\label{sec:repro}

Randomized search output is usually taken on trust; none of the numbers
above is.  Every reported system is discharged as data: the winning
state's chains are replayed deterministically (identical splitmix64 RNG
on GPU and CPU), each chain is shortened by the greedy reduction of
\cite[Algorithm 2]{OLGB21}, and the system is asserted to connect the
two projections through pairwise $1$-contiguous simplicial maps over
every facet of every domain.  An independent auditor---a pure-Python
program that rebuilds each triangulation from scratch and shares no
code with the search---re-runs every check; it has caught real bugs
during development and exits nonzero on any discrepancy.

The reproduction is a single script of audited, successful executions:
explicit product certificates for the circle, $S^2$, $S^3$, and $S^4$,
each re-verified and re-audited, complete in about two minutes on CPU
alone; an optional GPU mode replays the flagship annealer runs ($22$\,s,
$153$\,s, and $13$\,min).  Certificate documents listing every domain's
facets and every planner map are generated from data by a single
generator program; none of the tables in this paper is hand-typed from
anything but audited output.

\paragraph{Repositories.}  The project is split across two public repositories.
The \emph{source code}---the chain-first annealer, the certificate builder,
the independent auditor, and the reproduction script---lives at
\begin{center}
\url{https://github.com/MiguelPerezI/chain-first-annealer}
\end{center}
The \emph{certificates}---every domain and every planner map of every
certified system reported here (as \texttt{cover.txt}, \texttt{planners.txt}
and \texttt{summary.json} per run, and as readable certificate documents)---are
archived separately at
\begin{center}
\url{https://github.com/MiguelPerezI/simplicial-complexity-certificates}
\end{center}
A reader who wishes to verify, say, the $4$-sphere cover by hand can
open \texttt{cover.txt} and \texttt{planners.txt} in
\texttt{results/s4\_k3\_v2/} of the certificates repository, or run the
auditor (from the source repository) directly:\\
\texttt{python3 independent\_check.py results/s4\_k3\_v2}.

\section{Concluding remarks}\label{sec:concl}

The arc of this project is that the expensive object turned out to be
the wrong object.  The randomized greedy hunts chains one facet at a
time, at a cost that defeated a $96$-core cluster; the corrected
annealer hunts whole systems of chains and solves the classical targets
in seconds to minutes; and the rectangles those searches kept
discovering are a theorem that makes the search itself optional for a
whole family of spaces.  What remains genuinely open---exact values on
odd-dimensional spheres---now has the right shape: a homological filter
for candidate domains and an instrument that finds rectangles whenever
rectangles are what there is to find.

\begin{conjecture}\label{conj}
\emph{(C1)} $\SCst(\partial\Delta^4)=2$ at level $0$; equivalently,
with the convention $\SCst=(\text{minimum number of domains})-1$ used
throughout, two domains do not suffice on the S$^3$ model.  We have no
impossibility proof, and a failed search is not one: the support is
experimental.  Twelve two-chain searches (old and corrected energy, all
seeds) stall at coverage $460/500$ facets, and the largest single
two-chain domain we have observed has $329$ facets.
\emph{(C2)} On $\partial\Delta^{n+1}\times\partial\Delta^{n+1}$, greedy
and annealed searches with $k<\rho$ chains always converge to
(sub)rectangle covers of $\rho-1$ rectangles---the observed plateaus
$84/96$, $460/500$ and $2380/2520$ are the general case.
\end{conjecture}

\begin{thebibliography}{10}

\bibitem{BPV23}
A.~Borat, M.~Pamuk, and T.~Vergili.
\newblock Contiguity distance between simplicial maps.
\newblock {\em Turkish J. Math.}, 47:664--677, 2023.

\bibitem{Cerny85}
V.~\v{C}ern\'y.
\newblock Thermodynamical approach to the traveling salesman problem:
  an efficient simulation algorithm.
\newblock {\em J. Optim. Theory Appl.}, 45(1):41--51, 1985.

\bibitem{F03}
M.~Farber.
\newblock Topological complexity of motion planning.
\newblock {\em Discrete Comput. Geom.}, 29(2):211--221, 2003.

\bibitem{FTMV18}
D.~Fern\'andez-Ternero, E.~Mac\'ias-Virg\'os, E.~Minuz, and J.~A.
  Vilches.
\newblock Discrete topological complexity.
\newblock {\em Proc. Amer. Math. Soc.}, 146(10):4535--4548, 2018.

\bibitem{G18}
J.~Gonz\'alez.
\newblock Simplicial complexity: piecewise linear motion planning in
  robotics.
\newblock {\em New York J. Math.}, 24:279--292, 2018.

\bibitem{HD}
E.~Mac\'ias-Virg\'os and D.~Mosquera-Lois.
\newblock Homotopic distance between maps.
\newblock {\em Math. Proc. Cambridge Philos. Soc.}, 172(1):73--93, 2022.

\bibitem{KGV83}
S.~Kirkpatrick, C.~D. Gelatt, Jr., and M.~P. Vecchi.
\newblock Optimization by simulated annealing.
\newblock {\em Science}, 220(4598):671--680, 1983.

\bibitem{Metropolis53}
N.~Metropolis, A.~W. Rosenbluth, M.~N. Rosenbluth, A.~H. Teller, and
  E.~Teller.
\newblock Equation of state calculations by fast computing machines.
\newblock {\em J. Chem. Phys.}, 21(6):1087--1092, 1953.

\bibitem{Numba}
S.K. Lam, A.~Pitrou, and S.~Seibert.
\newblock Numba: a LLVM-based compiling Python JIT compiler.
\newblock In {\em Proc. Second Workshop on the LLVM Compiler
  Infrastructure in HPC}, article 7, 2015.

\bibitem{OLGB21}
C.~Ortiz, A.~Lara, J.~Gonz\'alez, and A.~Borat.
\newblock A randomized greedy algorithm for piecewise linear motion
  planning.
\newblock {\em Mathematics}, 9(19):2358, 2021.
\newblock arXiv:2008.13290.

\bibitem{Spanier}
E.H. Spanier.
\newblock {\em Algebraic topology}.
\newblock McGraw--Hill, New York, 1966.

\end{thebibliography}

\end{document}



